\chapter{Architecture logicielle}
\label{ch:arch}

\section{Vue en couches}

L'application suit une architecture en \textbf{quatre projets} alignés sur une \emph{clean architecture} simplifiée :

\begin{enumerate}[leftmargin=*]
  \item \textbf{Core} : aucune dépendance vers Infrastructure ni Web ; contient le modèle conceptuel (entités, enums), les contrats (\texttt{interface}) et les DTO utilisés par l'UI ou les services.
  \item \textbf{Infrastructure} : dépend de Core ; implémente la persistance (EF Core, SQLite), le \texttt{DbContext} et les migrations.
  \item \textbf{Services} : dépend de Core et Infrastructure ; concentre la logique métier, le calcul du portefeuille, les règles de transaction, l'analytique et l'intégration HTTP pour les cours.
  \item \textbf{Web} : point d'entrée ASP.NET Core ; configure l'injection de dépendances, le pipeline HTTP, Blazor et référence les trois autres couches.
\end{enumerate}

\section{Principes appliqués}

\begin{description}
  \item[Dépendances vers l'intérieur] L'UI et le composition root (\texttt{Program.cs}) résolvent les interfaces vers des implémentations concrètes enregistrées dans le conteneur IoC.
  \item[Testabilité] Les services prennent en charge des abstractions (\texttt{DbContext}, \texttt{HttpClient} typé, \texttt{IMemoryCache}, \texttt{IOptions}) injectables.
  \item[Cohésion] Un fichier \texttt{ServiceCollectionExtensions} centralise l'enregistrement des services métier et la configuration des options (\texttt{Alerts}, \texttt{PriceData}).
\end{description}

\section{Flux typique d'une requête utilisateur}

Pour une action sur le tableau de bord :
\begin{enumerate}[leftmargin=*]
  \item le composant Blazor (mode serveur) appelle un service injecté (\texttt{IPortfolioService}, \texttt{IAnalyticsService}, \texttt{IPortfolioAlertService}) ;
  \item le service interroge \texttt{AppDbContext} de manière asynchrone ;
  \item les résultats sont projetés en modèles (\texttt{PortfolioSnapshot}, \texttt{AssetAnalytics}, \texttt{PortfolioAlert}) ;
  \item l'UI met à jour l'état et, le cas échéant, invoque du JavaScript (Chart.js) pour le rendu graphique.
\end{enumerate}

\section{Diagramme de classes}

Un diagramme de classes Mermaid est fourni dans le dépôt à la racine : \texttt{DiagrammeClasses.md}. Il est recommandé d'\textbf{exporter} les figures (PNG/SVG) depuis \url{https://mermaid.live} et de les insérer ici avec :

\begin{verbatim}
\begin{figure}[h]
  \centering
  \includegraphics[width=\textwidth]{figures/domaine-ef.png}
  \caption{Modèle entités-associations (extrait).}
\end{figure}
\end{verbatim}

Créer un dossier \texttt{Rapport/figures/} et y placer les exports pour que la compilation \LaTeX{} réussisse.
